% Exemplo: modelo de apresentação acadêmica (tema metropolis, 16:9) com colunas, figura,
% tabela, citações com biblatex-abnt, referências e slide de encerramento
% Compilar: latexmk -pdf modelo-apresentacao.tex   Testado com: beamer, metropolis, biblatex-abnt, TeX Live 2024+
% Ponto de partida para defesas e seminários: troque os dados, o conteúdo e, se quiser, o logo
% da sua instituição (\titlegraphic). A imagem "example-image" vem do pacote mwe.
%%BEGIN:modelopreambulo
\documentclass[aspectratio=169]{beamer}
\usepackage[T1]{fontenc}
\usepackage[brazilian]{babel}
\usetheme[progressbar=frametitle]{metropolis}
\usepackage{booktabs}
\usepackage{csquotes}
\usepackage[backend=biber, style=abnt]{biblatex}
\addbibresource{../referencias/catalogo.bib}

\title{Título do trabalho}
\subtitle{Defesa de Trabalho de Conclusão de Curso}
\author{Nome do Autor \\ \small Orientador: Prof. Dr. Fulano de Tal}
\institute{Universidade Federal de Viçosa}
\date{\today}
% \titlegraphic{\hfill\includegraphics[height=1cm]{logo}}
%%END:modelopreambulo

\begin{document}

\maketitle

\begin{frame}{Roteiro}
  \tableofcontents
\end{frame}

\section{Introdução}

%%BEGIN:modelocolunas
\begin{frame}{Problema e objetivos}
  \begin{columns}[T]
    \begin{column}{.55\textwidth}
      \begin{itemize}
        \item Contexto do problema \parencite{ex-livro}
        \item Objetivo geral
        \item Objetivos específicos
      \end{itemize}
    \end{column}
    \begin{column}{.4\textwidth}
      \includegraphics[width=\linewidth]{example-image}
    \end{column}
  \end{columns}
\end{frame}
%%END:modelocolunas

\section{Resultados}

\begin{frame}{Resultados}
  \centering
  \begin{tabular}{lrr}
    \toprule
    Cenário & Acurácia (\%) & Tempo (s) \\
    \midrule
    A & 91,2 & 3,4 \\
    B & 94,8 & 5,1 \\
    \bottomrule
  \end{tabular}
\end{frame}

%%BEGIN:modelobib
\begin{frame}[allowframebreaks]{Referências}
  \printbibliography[heading=none]
\end{frame}
%%END:modelobib

\begin{frame}[standout]
  Obrigado!
\end{frame}

\end{document}
